% Generated draft: requires mathematical and source review.
\documentclass[11pt]{article}
\usepackage[margin=25mm]{geometry}
\usepackage{fontspec}
\setmainfont{Latin Modern Roman}
\usepackage{amsmath,amssymb,mathtools,unicode-math}
\setmathfont{Latin Modern Math}
\usepackage{enumitem,needspace,xurl,hyperref}
\setlength{\parindent}{0pt}
\setlength{\parskip}{5pt}
\setlength{\emergencystretch}{3em}
\newcommand{\N}{\mathbb{N}}
\newcommand{\Z}{\mathbb{Z}}
\newcommand{\Q}{\mathbb{Q}}
\newcommand{\R}{\mathbb{R}}
\newcommand{\C}{\mathbb{C}}
\newcommand{\F}{\mathbb{F}}
\newcommand{\A}{\mathbb{A}}
\newcommand{\PP}{\mathbb{P}}
\newcommand{\E}{\mathbb{E}}
\newcommand{\eps}{\varepsilon}
\newcommand{\dd}{\,\mathrm{d}}
\begin{document}
\section*{UnsolvedMath ID: 700874}
\textbf{Source problem code:} catalogue entry
\subsection{Definitions and mathematical statement}
For a prime number \(p\) and an integer \(n\geq 5\), let
\[
f(n,p)=\#\{[G]_{\cong}:G\text{ is a finite group and }|G|=p^n\},
\]
where \([G]_{\cong}\) denotes the isomorphism class of \(G\). Counting groups up to isomorphism is the conventional interpretation of the supplied wording; the source record does not explicitly specify this convention.

The question attributed to A. Mann asks whether, for every fixed integer \(n\geq 5\), the function \(p\mapsto f(n,p)\), with domain the prime numbers, is increasing. The supplied record does not specify whether increasing means strictly increasing or nondecreasing. Thus the two possible readings are
\[
\forall n\in\Z,\ n\geq 5,\quad
\forall p,q\text{ prime},\quad
p<q\ \Longrightarrow\ f(n,p)<f(n,q),
\]
or, under the nondecreasing interpretation,
\[
\forall n\in\Z,\ n\geq 5,\quad
\forall p,q\text{ prime},\quad
p<q\ \Longrightarrow\ f(n,p)\leq f(n,q).
\]
The intended inequality requires clarification from the source.
\subsection{Short English statement}
For each fixed integer \(n\geq 5\), does the number of isomorphism classes of groups of order \(p^n\) increase as the prime \(p\) increases? The supplied wording leaves strict versus weak increase unspecified.
\subsection{Sources}
\begin{itemize}
\item Supplied local UnsolvedMath registry record: permanent ID \(700874\), source code \texttt{catalogue entry}, title \emph{Kourovka Notebook Problem 12.55}, upstream revision \texttt{d371ee917ac9647fbb06cd6fc529deec6f3dda82}.
\item A. Mann, Problem 12.55, \emph{The Kourovka Notebook (Unsolved Problems in Group Theory)}, 12th issue (1992), as attributed in the supplied record. The record also identifies the 21st issue edition (2026), \texttt{arXiv:1401.0300v48}. The supplied collection description names E. I. Khukhro and V. D. Mazurov as editors.
\item The registry labels this problem \emph{open}. This is a supplied status label, not an independently verified mathematical or literature assessment. The cited source text has not been independently checked for this draft.
\end{itemize}
\end{document}
